\section{Discrepancy bounds: further entries}\label{s:kimb}
In 2017 Kimberling added to several hundred OEIS entries on binary morphic words a conjecture
$\ell<nr-a(n)<h$, where $a(n)$ is the position of the $n$-th occurrence of a letter and $1/r$ is its
frequency. Part~\ref{part:one} treats the nine entries whose morphisms appear in our original note. Here the
same method (Section~\ref{s:method}, script \texttt{python/discrepancy\_exact.py}, unchanged) is applied to
all other entries of this kind that we found and that were not already settled. For each entry the positions
computed from the morphism are compared with the data of the entry, and the constant $r$ stated by
Kimberling is compared with $1/f_L$ computed exactly from the Perron eigenvector.

When an entry speaks of a ``$c$-limiting word'' of $\sigma$ (because $\sigma(c)$ does not begin with $c$),
the word is the fixed point of $\sigma^2$ beginning with $c$, and the method is applied to $\sigma^2$.

\paragraph{Attainment.} In every case below $r$ is irrational, so $D(n)=nr-a(n)$ is irrational for $n\ge1$ and
is never equal to an integer bound. When a conjectured integer bound equals the exact infimum or supremum,
the strict inequality is therefore true and the bound is optimal. A non-integer bound equal to $r$ cannot be
attained either, since $D(n)=r$ would give $(n-1)r=a(n)$.

\subsection{First series: fixed points}
\begin{theorem}\label{t:kimb1} For the following entries the exact extreme values of $D(n)$ are as
stated, and Kimberling's conjecture is true.\end{theorem}
\begin{center}\small\begin{tabular}{llllll}\toprule
OEIS & morphism, letter & $r$ & $\inf D$ & $\sup D$ & conjecture\\\midrule
\oeis{A026363} & $0\to11$, $1\to101$, 1 & $(1+\sqrt3)/2$ & $-1$ & $(1+\sqrt3)/2$ & $(-1,2)$\\
\oeis{A026364} & $0\to11$, $1\to101$, 0 & $2+\sqrt3$ & $-1$ & $2+\sqrt3$ & $(-1,4)$\\
\oeis{A026368} & $0\to11$, $1\to110$, 0 & $2+\sqrt3$ & $-1-\sqrt3/3$ & $2+2\sqrt3/3$ & $(-2,4)$\\
\oeis{A284753} & $0\to01$, $1\to0001$, 1 & $1+\sqrt3$ & $-2$ & $1+\sqrt3$ & $(-2,3)$\\
\oeis{A284894} & $0\to01$, $1\to0111$, 0 & $2+\sqrt2$ & $1+\sqrt2$ (at $n=1$) & $3+2\sqrt2$ & $(2,6)$\\
\oeis{A285669} & $0\to11$, $1\to1101$, 0 & $(5+\sqrt{17})/2$ & $(-7+\sqrt{17})/16$ & $(29+5\sqrt{17})/16$ & $(-1,4)$\\\bottomrule
\end{tabular}\end{center}
For \oeis{A284753}, \oeis{A284894} and \oeis{A285669} only boundedness was known (Dekking, from
Adamczewski's theorem~\cite{Ad}). Dekking notes in \oeis{A284753} that it would imply
$-1<n(1+\sqrt3)/2-\oeis{A026363}(n)<3/2$; this stronger statement is true, since $\sup=(1+\sqrt3)/2<3/2$.

\subsection{Second series: limiting words}
\begin{theorem}\label{t:kimb2} For the following entries Kimberling's conjecture is true, with the exact
extreme values shown.\end{theorem}
\begin{center}\small\begin{tabular}{llllll}\toprule
OEIS & $\sigma$, start, letter & $\inf D$ & $\sup D$ & conjecture\\\midrule
\oeis{A026367} & $0\to11$, $1\to110$; letter 1 & $-(3+\sqrt3)/6$ & $1+\sqrt3/3$ & $(-1,2)$\\
\oeis{A284654} & $0\to1$, $1\to0110$; 0-limit, letter 0 & $-1$ & $(3+\sqrt3)/2$ & $(-1,3)$\\
\oeis{A284675} & $0\to1$, $1\to0111$; 0-limit, letter 0 & $(1+\sqrt{13})/2$ & $(3+\sqrt{13})/2$ & $(2,4)$\\
\oeis{A284676} & same, letter 1 & $(-5+\sqrt{13})/2$ & $(-3+\sqrt{13})/2$ & $(-1,4)$\\
\oeis{A284678} & $0\to1$, $1\to0111$; 1-limit, letter 0 & $(1+\sqrt{13})/2$ & $(3+\sqrt{13})/2$ & $(2,4)$\\
\oeis{A284679} & same, letter 1 & $(-5+\sqrt{13})/2$ & $(-3+\sqrt{13})/2$ & $(-1,1)$\\
\oeis{A285032} & $0\to10$, $1\to001$; 0-limit, letter 0 & $-1+\sqrt2/4$ & $1/2+3\sqrt2/4$ & $(-1,2)$\\
\oeis{A285035} & $0\to10$, $1\to001$; 1-limit, letter 0 & $-1+\sqrt2/4$ & $1/2+3\sqrt2/4$ & $(-1,2)$\\
\oeis{A285036} & same, letter 1 & $-3/2+\sqrt2/2$ & $1/2+\sqrt2$ & $(-1,2)$\\
\oeis{A285082} & $0\to10$, $1\to011$; 0-limit, letter 1 & $-1$ & $1$ & $(-1,1)$\\
\oeis{A285084} & $0\to10$, $1\to011$; 1-limit, letter 0 & $0$ & $r=(3+\sqrt5)/2$ & $(0,r)$\\
\oeis{A285275} & $0\to10$, $1\to0111$; 0-limit, letter 0 & $0$ & $3+2\sqrt2$ & $(0,6)$\\
\oeis{A285278} & $0\to10$, $1\to0111$; 1-limit, letter 0 & $0$ & $3+2\sqrt2$ & $(0,6)$\\
\oeis{A285419} & $0\to11$, $1\to011$; 0-limit, letter 0 & $-1+\sqrt3/3$ & $2+4\sqrt3/3$ & $(-1,5)$\\\bottomrule
\end{tabular}\end{center}
In \oeis{A285082} and \oeis{A285084} both conjectured bounds are the exact infimum and supremum, and
neither is attained.

\subsection{Third series}
The third series contains the companion entries of the words of Part~\ref{part:one} (\oeis{A283964},
\oeis{A283967}, \oeis{A284370}, \oeis{A284506}, \oeis{A285359}) and further families. In the table, $\tau$ is
the morphism to which the method is applied ($\sigma^2$ for a limiting word), with the first letter of the
fixed point in parentheses, followed by the letter whose positions are listed. All the values of $r$ stated
in these entries agree with $1/f_L$, except in \oeis{A285279} (Section~\ref{s:kimbfalse}).
\begin{theorem}\label{t:kimb3} For each entry in the following table, Kimberling's conjecture is true,
and the exact infimum and supremum of $D(n)$ are the values shown.\end{theorem}
{\footnotesize\setlength{\tabcolsep}{4pt}\begin{longtable}{llllll}\toprule
OEIS & $\tau$ (start), letter & $r$ & $\inf D$ & $\sup D$ & conj.\\\midrule\endhead
\tinput tables/kimb3.tex 
\oeis{A284015} & $0\to1,1\to10101$ (1), 1 & $\tfrac{-1+\sqrt{17}}{2}$ & $\tfrac{-7+\sqrt{17}}{4}$ & $\tfrac{1+\sqrt{17}}{4}$ & $(-1,2)$\\\bottomrule
\end{longtable}}

\subsection{False conjectures}\label{s:kimbfalse}
\begin{theorem}\label{t:kimbfalse} The conjectures of the following entries are false as stated. The value
of $n$ is the first one for which $D(n)$ lies outside the conjectured interval, computed in exact arithmetic
with the constant $r$ stated in the entry; in each case $a(n)$ is in the data of the entry.\end{theorem}
\begin{center}\small\begin{tabular}{llllll}\toprule
OEIS & stated $r$ & conjecture & $n$ & $a(n)$ & $D(n)$\\\midrule
\oeis{A285129} & $1.3903\ldots$ & $(-1,1)$ & 2 & 4 & $-1.219\ldots$\\
\oeis{A285130} & $3.5621\ldots$ & $(-1,3)$ & 2 & 3 & $4.124\ldots$\\
\oeis{A285134} & $1+\sqrt{1/3}$ & $(-1,1)$ & 2 & 2 & $2\sqrt3/3$\\
\oeis{A285135} & $1+\sqrt3$ & $(-1,1)$ & 1 & 4 & $-3+\sqrt3$\\
\oeis{A285137} & $1+\sqrt{1/3}$ & $(-1,1)$ & 3 & 6 & $-3+\sqrt3$\\
\oeis{A285138} & $1+\sqrt3$ & $(-1,\sqrt3)$ & 2 & 3 & $-1+2\sqrt3$\\
\oeis{A285140} & $(3+\sqrt3)/3$ & $(-1,1)$ & 2 & 2 & $2\sqrt3/3$\\
\oeis{A285141} & $1+\sqrt3$ & $(-1,1)$ & 6 & 15 & $-9+6\sqrt3$\\
\oeis{A285143} & $(3+\sqrt3)/3$ & $(-1,1)$ & 12 & 20 & $-8+4\sqrt3$\\
\oeis{A285144} & $1+\sqrt3$ & $(-1,\sqrt3)$ & 2 & 3 & $-1+2\sqrt3$\\
\oeis{A285279} & $2+\sqrt2$ & $(0,6)$ & 4 & 6 & $2+4\sqrt2$\\\bottomrule
\end{tabular}\end{center}
For \oeis{A285138} and \oeis{A285144}, $D(1)=\sqrt3$ is equal to the bound, and the first strict violation is
at $n=2$. For \oeis{A285279} the stated $r$ is not the inverse frequency, which is $\sqrt2$. Each
counterexample was re-checked on November 19, 2026 against the data of the entry, with the constant as
printed there. \oeis{A285129} and \oeis{A285130} come from the non-Pisot morphism $0\mapsto10$,
$1\mapsto0000$ of Section~\ref{s:A285126} (its $1$-limiting word).

\begin{proposition} The conjecture $-1+\sqrt2<n(1+\sqrt2)-a(n)<\sqrt2$ of \oeis{A285078} is false for $n=1$,
where $a(1)=1$ and $D(1)=\sqrt2$, and true for all $n\ge2$; both bounds are the exact infimum and supremum
over $n\ge2$.\end{proposition}
This is the same phenomenon as the one noted by Dekking for \oeis{A285074}.

\subsection{A non-Pisot morphism}\label{s:A285126}
\begin{proposition} Let $\sigma:0\mapsto10$, $1\mapsto0000$, and let \oeis{A285125} be its $0$-limiting word
(the fixed point of $\sigma^2$ beginning with $0$). The conjectures $0<nr-a(n)<1$ of \oeis{A285126} and
$-3<ns-b(n)<0$ of \oeis{A285127}, with $r=1.39038\ldots$ and $s=3.56155\ldots$ the inverse frequencies,
are false: $3r-a(3)=3r-3=1.171\ldots$ and $2s-b(2)=2s-11=-3.876\ldots$. Moreover the discrepancy is
unbounded.\end{proposition}
\begin{proof} The counterexamples are in the data. The incidence matrix has eigenvalues $(1\pm\sqrt{17})/2$,
and the second one, $-1.56\ldots$, has modulus larger than~$1$. The Perron eigenvector is
$(4,\lambda-1)$ with $\lambda=(1+\sqrt{17})/2$, which gives the frequencies. By the representation
\eqref{e:D} with $|\lambda'|>1$, $D$ is unbounded; numerically, on $3\cdot10^6$ letters, $nr-a(n)$ reaches
$1017.3$ and $ns-b(n)$ reaches $-2603.3$. \end{proof}

\subsection{Transforms of the Fibonacci word}\label{s:fibtr}
Let $F=\oeis{A003849}$ be the Fibonacci word, fixed point of $0\mapsto01$, $1\mapsto0$, and let
$T_1:0\mapsto1,1\mapsto001$, $T_2:0\mapsto1,1\mapsto010$, $T_3:0\mapsto1,1\mapsto011$.
\begin{theorem}\label{t:fibtr} The following conjectures of Kimberling are true:
\begin{itemize}
\item \oeis{A287676} (positions of $1$ in $T_1(F)$): $-1<n(4-\sqrt5)-a(n)<2$; moreover Liu's formula
  $a(n)=5n+2-2\lfloor(n+1)\varphi\rfloor$ holds, so $n(4-\sqrt5)-a(n)=\sqrt5-1-2\{(n+1)\varphi\}$, with exact
  bounds $\sqrt5-3$ and $\sqrt5-1$;
\item \oeis{A287724} (positions of $1$ in $T_2(F)$): $0<n(4-\sqrt5)-a(n)<3$;
\item \oeis{A287723} (positions of $0$ in $T_2(F)$): $-1<n(7+\sqrt5)/4-a(n)<1$;
\item \oeis{A287727} (positions of $1$ in $T_3(F)$): $-1<n(15-\sqrt5)/10-a(n)<1$.
\end{itemize}\end{theorem}
\begin{proof}[Proof (\Walnut, Fibonacci numeration; \texttt{fibtransform/fib.wal})] The map
$n\mapsto\lfloor n\varphi\rfloor$ is defined through the Zeckendorf shift, and \Walnut{} checks that it is
strictly increasing and that $\{\lfloor n\varphi\rfloor\}$ and $\{\lfloor n\varphi\rfloor+n\}$ partition the
positive integers, which characterises it (Lambek--Moser). The number of $1$s in $F[0..i)$ is
$2i+1-\lfloor(i+1)\varphi\rfloor$ (checked by induction in \Walnut), so the block $T(F[i])$ starts at
$1+i+2\cdot\mathrm{ones}(i)$, and the positions of the letters follow block by block; these are checked to be
functions of $n$ and to agree with the data. The constants are in $\mathbb{Q}(\varphi)$:
$4-\sqrt5=5-2\varphi$, $(7+\sqrt5)/4=(\varphi+3)/2$, $(15-\sqrt5)/10=(8-\varphi)/5$, and each inequality
becomes an integer comparison with $\lfloor n\varphi\rfloor$ or $\lfloor2n\varphi\rfloor$, because
$n\varphi$ is irrational. All four statements and Liu's formula are \TRUE; the control with a tightened bound
is \FALSE. \end{proof}

\section{Limit conjectures \texorpdfstring{$a(n)/n\to m$}{a(n)/n -> m}}\label{s:limits}
\begin{lemma}[Queffélec~\cite{Qu}]\label{l:perron} Let $\sigma$ be a primitive substitution and $x$ a fixed
point of a power of $\sigma$. Every letter $L$ has a frequency $f_L$ in $x$, equal to the $L$-component of the
Perron eigenvector of the incidence matrix normalised to sum $1$. Hence the position $a(n)$ of the $n$-th
$L$ satisfies $a(n)/n\to1/f_L$. More generally every factor $u$ has a frequency, given by the Perron
eigenvector of the substitution induced on blocks of length $|u|$.\end{lemma}

\subsection{Morphic words}
For each entry below the morphism is primitive ($M^2$ or $M^3$ has positive entries, checked), the positions
recomputed from the morphism agree with the data, and $1/f_L$ is computed exactly and compared symbolically
with the stated limit. For transforms $\{00\to0,11\to1\}$ of a word made of aligned pairs, the frequency is the
one of the base word.
\begin{theorem}\label{t:limits} The limits stated by Kimberling are true for the entries marked ``true'' and
false for those marked ``false'' (the limit exists but has another value).\end{theorem}
{\small\begin{longtable}{lllll}\toprule
OEIS & exact limit $1/f_L$ & stated limit & observed & result\\\midrule\endhead
\tinput tables/limits.tex 
\bottomrule\end{longtable}}
The column ``observed'' is $a(N)/N$ at the end of the data computed (about $2.5\cdot10^5$ terms). In the false
cases the stated value is the limit of another entry of the same family (for instance $(7+\sqrt{13})/6$ for a
word whose eigenvalue is in $\mathbb{Q}(\sqrt3)$), which suggests a copying slip.

\subsection{Transforms of the Thue--Morse word}\label{s:tm}
For a word $u$ and a letter $c$, $\{u\to c\}(\mathrm{TM})$ replaces each occurrence of $u$ in the Thue--Morse
word by $c$, from left to right. The occurrences used ($u\in\{00,11,01,10,011,101,0011,1001,1100\}$) cannot
overlap, since an overlap would create a factor $axaxa$, and the Thue--Morse word is overlap-free. If $f_u$
is the frequency of $u$, the frequency of $L$ in the transform is
\[ \frac{1/2-|u|_L f_u+[c=L]f_u}{1-(|u|-1)f_u}, \]
and the $f_u$ are computed exactly by Lemma~\ref{l:perron} (eigenvalue $2$), then checked empirically.
\begin{theorem} The limits stated by Kimberling are true for the following entries, except \oeis{A285964},
where the stated value $4/5$ is impossible ($a(n)\ge n$) and the true limit is $5/4$.\end{theorem}
\begin{center}\small\begin{longtable}{lllll}\toprule
OEIS & transform, letter & exact limit & stated & result\\\midrule\endhead
\tinput tables/tm.tex 
\bottomrule\end{longtable}\end{center}
For \oeis{A003157}, the identification with the positions of $1$ in \oeis{A286044} is the one given in the
entry, checked on its data.

\subsection{Non-aligned transforms}\label{s:runs}
The greedy transform $\{00\to0,11\to1\}$ applied to the limiting words \oeis{A285565} and \oeis{A285568} of
$0\mapsto11$, $1\mapsto0110$ maps each maximal run of $k$ equal letters to $\lceil k/2\rceil$ letters (checked on
the data). The run frequencies are the frequencies of the factors $yx^ky$, computed exactly from the Perron
eigenvector ($\lambda=1+\sqrt5$) of the substitution induced on blocks of length (maximal run $+2$).
\begin{theorem} The limits $1+\sqrt5/2$ (letter $0$) and $1+2/\sqrt5$ (letter $1$) conjectured in
\oeis{A285590}, \oeis{A285591}, \oeis{A285593} and \oeis{A285594} are true.\end{theorem}

\subsection{Six limits equal to \texorpdfstring{$\varphi$}{phi}}\label{s:A285503}
\begin{theorem} In \oeis{A285503}, \oeis{A285517}, \oeis{A285567}, \oeis{A285623}, \oeis{A285627} and
\oeis{A285663}, $a(n)/n\to\varphi$.\end{theorem}
\begin{proof} The words are limiting words or fixed points of $0\mapsto11$, $1\mapsto w$ with
$w\in\{0011,0110,1001,1010,1100\}$; \oeis{A285517} reads \oeis{A285501} by pairs. All five morphisms have
incidence matrix $M=\begin{pmatrix}0&2\\2&2\end{pmatrix}$, with $M^2>0$, Perron eigenvalue $1+\sqrt5$ and
eigenvector $(2,1+\sqrt5)$. The frequency of $1$ is $(1+\sqrt5)/(3+\sqrt5)=1/\varphi$, and
Lemma~\ref{l:perron} applies. For \oeis{A285517}, all images have even length and are made of the blocks $00$
and $11$, so \oeis{A285501} is a sequence of aligned pairs and the pair word has the same frequency of
$1$. \end{proof}
The second eigenvalue $1-\sqrt5$ has modulus larger than $1$, so these words are not Pisot and
$a(n)-n\varphi$ is unbounded: only the limit is true, not any bound on the discrepancy.
